% Created 2025-10-06 Mon 22:30
% Intended LaTeX compiler: pdflatex
\documentclass[11pt]{article}
\usepackage[utf8]{inputenc}
\usepackage[T1]{fontenc}
\usepackage{graphicx}
\usepackage{longtable}
\usepackage{wrapfig}
\usepackage{rotating}
\usepackage[normalem]{ulem}
\usepackage{amsmath}
\usepackage{amssymb}
\usepackage{capt-of}
\usepackage{hyperref}
\author{gnalens}
\date{\textit{{[}2025-08-26 Tue 10:40]}}
\title{Deriving Relational Dynamics}
\hypersetup{
 pdfauthor={gnalens},
 pdftitle={Deriving Relational Dynamics},
 pdfkeywords={},
 pdfsubject={},
 pdfcreator={},
 pdflang={English}}
\usepackage{biblatex}
\addbibresource{~/ems/org/ref/bibliography.bib}
\begin{document}

\maketitle
This \href{derivation.org}{derivation} explores the consequences of the \href{claim.org}{claim} that some \href{thing.org}{thing} is \href{existing.org}{existing}.
\section*{Deriving Existential Coherence}
\label{sec:org3ae291d}
\begin{itemize}
\item \href{claiming.org}{Claiming} that any \href{thing.org}{thing} \href{existing.org}{exists} is not trivial.
\begin{itemize}
\item It is to claim:
\begin{itemize}
\item Any thing is different no thing, or at a minimum
\item Any thing is not what it is not.
\end{itemize}
\item That is: \texttt{thing not-equal not-thing}
\end{itemize}
\item This claim has \href{dynamic.org}{dynamics} which must be \href{deriving.org}{derived}:
\begin{itemize}
\item What is a \texttt{thing}, \texttt{not} (\href{negating.org}{negating}), and \texttt{equal} (\href{equaling.org}{equaling})?
\item How do they result in a thing \texttt{existing}?
\end{itemize}
\item Given a thing is \href{determining.org}{determined} by its \href{inequivalent.org}{inequivalence} to any thing that is not itself, it cannot not derive there is a \href{dynamic-of-determination.org}{dynamic of determination}.
\begin{itemize}
\item In the context of a thing existing, this dynamic is the \href{act.org}{acts} of \href{including.org}{including} and \href{excluding.org}{excluding}:
\item A thing includes itself it what it is.
\item That thing excludes every thing it is not.
\end{itemize}
\item Given that \href{negation.org}{negation} is determined by how it relates to a thing and \href{equivalence.org}{equivalence}, it cannot not derive there is a \href{dynamic-of-relation.org}{dynamic of relation}.
\begin{itemize}
\item In the context of a thing existing, this dynamic is the \href{condition.org}{conditions} of \href{inclusion.org}{inclusion} and \href{exclusion.org}{exclusion}:
\begin{itemize}
\item Inclusion is the condition of existing as a thing.
\item Exclusion is the condition of not existing that thing.
\end{itemize}
\end{itemize}
\item Given that equivalence is determined by how it relates a thing and not that thing, and is related to negation, it cannot not derive that there is a \href{dynamic-of-structuration.org}{dynamic of structuration}:
\begin{itemize}
\item In the context of a thing existing, this dynamic is the \href{structure.org}{structure} of \href{coherence.org}{coherence}:
\begin{itemize}
\item A \href{elementary-structure.org}{elementary structure} that \href{stabilizing.org}{stabilizes} the determination of relation across \href{tension.org}{tension} between inclusion and exclusion
\begin{itemize}
\item This coherence \href{inducing.org}{induces} \href{type-homogeneity.org}{type-homogeneity} in existing: \href{including.org}{including} and \href{excluding.org}{excluding} act on the same condition on the same structure.
\end{itemize}
\end{itemize}
\end{itemize}
\item Determining thing, negation, and equivalence provides a \href{dynamic-of-existing.org}{dynamic of existing}, but leaves undetermined dynamics which must be derived:
\begin{itemize}
\item How is the \href{elementary-structure-of-relational-coherence.org}{elementary structure of relational coherence} \href{stabilizing.org}{stabilized}?
\end{itemize}
\end{itemize}
\section*{Deriving Existential Stabilization}
\label{sec:org295c04d}
\begin{itemize}
\item Given that the \href{stable.org}{stable} coherence of existing is determined by the relation between inclusion and exclusion, it cannot not derive there is a dynamic of determination:
\begin{itemize}
\item In the context of a thing existing, this dynamic is the act of \href{relating.org}{relating}: the act that stabilizes the tension between including and excluding.
\item Given that the only undetermined dynamic is the stabilization of the elementary relational unit, relating must be the act of that unit determining its own derviation.
\item This relating induces \href{reflexive-sequencing.org}{reflexive sequencing}: a unit's existence cannot be determined without the unit's ongoing derivation of its existence. (Look, mathematics and metaphysics is hard, and language is also hard.)
\item 
\end{itemize}
\end{itemize}

\#\# Deriving Relational Coherence
\begin{itemize}
\item Given that the sustained coherence of existing is determined by the relation between inclusion and exclusion, it cannot not derive there is a dynamic of determination:
\begin{itemize}
\item In the context of something existing, this dynamic is the act of relating: the continuous act that maintains the interplay between including and excluding.
\item This act introduces reflexive sequence: each relational unit can relate to itself, producing an indefinitely extensible sequence of deepened relational units.
\item Reflexive lineage is the traceable sequence of relational structures generated by successive self-relation, preserving the ancestry of a unit across reflexive deepening.
\end{itemize}
\item Given that the act of relating is determined by how it engages and sustains inclusion and exclusion within a unified coherence, it cannot not derive there is a dynamic of relation:
\begin{itemize}
\item In the context of something existing, this dynamic is the condition of relation: the stabilized condition in which inclusion and exclusion are coherently maintained.
\item This condition induces positional invariance: inclusion and exclusion must always be distinct and ordered within any relational structure, preserving their roles through reflexive deepening.
\end{itemize}
\item Given that the condition of relation is determined by its reliance on the ongoing act of relating for its persistence, it cannot not derive there is a dynamic of structure:
\begin{itemize}
\item In the context of something existing, this dynamic is the structure of relational form: the minimal structural configuration that formalizes the interplay of relation and relating, ensuring the dynamic's coherence.
\item This structure establishes reflexive compositionality: relational forms at any depth can be composed side-by-side into ordered composites, preserving coherence across both depth and breadth.
\item A composite is an ordered configuration of two or more relational units, held together at a shared reflexive sequence to form a structured whole.
\end{itemize}
\item Determining relating, relation, and relational form provides the minimal dynamic determination for cohering the relational dynamic unit, but leaves unresolved dynamics which must be derived:
\begin{itemize}
\item What coheres relational form itself?
\end{itemize}
\end{itemize}
\#\# Deriving Sustained Relational Dynamics
\begin{itemize}
\item Given that relational form is determined by determining relating and relation, it cannot not derive there is a dynamic of determination:
\begin{itemize}
\item In the context of something existing, this is the act of relating-relation: the act that maintains and mediates between the act of relating and the condition of relation, ensuring their joint persistence within relational form.
\end{itemize}
\item Given that relating-relation is determined by how it connects and sustains relating and relation within relational form, it cannot not derive there is a dynamic of relation:
\begin{itemize}
\item In the context of something existing, this is the condition of relational self-coherence: the condition in which relating and relation are not only maintained together but are stabilized as a unified dynamic.
\end{itemize}

\item Given that relational self-coherence is determined by how it structures the persistent relation between relating and relation, it cannot not derive there is a dynamic of structure:

\begin{itemize}
\item In the context of something existing, this is the structure of relational closure: a structured configuration that formalizes the self-maintenance of relation through relating, aiming to achieve closure and internal consistency.
\end{itemize}

\item Determining relating-relation, relational self-coherence, and relational closure provides the minimal dynamic determination for cohering the self-sustaining relational unit, but leaves unresolved dynamics which must be determined:

\begin{itemize}
\item What determines the coherence of the self-sustaining relational unit?
\end{itemize}
\end{itemize}
\#\# Deriving Bounded Relational Coherence
\begin{itemize}
\item Given that the coherence of the self-sustaining relational unit is determined by its need to distinguish itself from what it is not, it cannot not derive there is a dynamic of determination:

\begin{itemize}
\item In the context of something existing, this is the act of boundary-excluding: the act that enacts the exclusion of what lies beyond the relational unit—relationlessness.
\end{itemize}

\item Given that the act of boundary-excluding is determined by how it engages with what stands outside the relational unit, it cannot not derive there is a dynamic of relation:

\begin{itemize}
\item In the context of something existing, this is the condition of relational negation: the condition that maintains the distinction between relation and what is excluded from relation.
\end{itemize}

\item Given that the condition of relational negation is determined by how it structures the distinction between relation and non-relation, it cannot not derive there is a dynamic of structure:

\begin{itemize}
\item In the context of something existing, this is the structure of relational boundary: the minimal structural form that formalizes and stabilizes the division between the relational unit and what is outside it.
\end{itemize}

\item Determining boundary-excluding, relational negation, and relational boundary provides the minimal dynamic determination for cohering the relational boundary unit, but leaves unresolved dynamics which must be derived:

\begin{itemize}
\item What determines the coherence of the relational boundary unit?
\end{itemize}
\end{itemize}
\#\# Deriving Reflexive Relational Coherence
\begin{itemize}
\item Given that the coherence of the relational boundary unit is determined by its distinction between relation and non-relation, it cannot not derive there is a dynamic of determination:

\begin{itemize}
\item In the context of something existing, this is the act of reflexive relating: the act that enacts relation upon the relational boundary itself, engaging the very distinction it establishes.
\end{itemize}

\item Given that the act of reflexive relating is determined by how it relates the relational boundary back to the relational field, it cannot not derive there is a dynamic of relation:

\begin{itemize}
\item In the context of something existing, this is the condition of reflexive relation: the condition in which the relational boundary becomes part of the relational system through its own self-engagement.
\end{itemize}

\item Given that the condition of reflexive relation is determined by how it structures the self-engagement of the relational boundary within relation, it cannot not derive there is a dynamic of structure:

\begin{itemize}
\item In the context of something existing, this is the structure of reflexive form: the minimal structure that formalizes the folding of the relational boundary into the relational system, ensuring that the act of exclusion is itself related and coherently stabilized.
\end{itemize}

\item Determining reflexive relating, reflexive relation, and reflexive form provides the minimal dynamic determination for cohering the reflexive relational unit, but leaves unresolved dynamics which must be derived:
\begin{itemize}
\item What coheres the reflexive relational unit?
\end{itemize}
\end{itemize}
\#\# Deriving Multirelational Coherence
\begin{itemize}
\item Given that the coherence of the reflexive relational unit is determined by its differentiation from what is not itself, it cannot not derive there is a dynamic of determination:

\begin{itemize}
\item In the context of something existing, this is the act of external differentiation: the act by which a reflexive relational unit distinguishes itself not only from non-relation but from other distinct reflexive relational units within the relational field.
\end{itemize}

\item Given that the act of external differentiation is determined by how it positions the reflexive relational unit relative to others, it cannot not derive there is a dynamic of relation:

\begin{itemize}
\item In the context of something existing, this is the condition of co-presence: the condition in which multiple reflexive relational units are maintained in differentiated yet shared existence within the relational field.
\item This condition enables relational tension: distinct units held in co-presence generate structured tension, which preserves both their differentiation and mutual relationality.
\item Operation types—such as Tension, Composition, and Alignment—specify the relational acts by which units are bound, ordered, or depth-aligned; their nature remains invariant even as structures deepen reflexively.
\end{itemize}
\item Given that the condition of co-presence is determined by how it sustains both differentiation and relational coherence between multiple units, it cannot not derive there is a dynamic of structure:

\begin{itemize}
\item In the context of something existing, this is the structure of relational tension: the minimal structure that formalizes and stabilizes the interplay between reflexive relational units, ensuring both their distinctness and their mutual relational coherence.
\item This structure supports ordered chain composition: multiple tension units can be assembled into sequenced composites, forming complex relational architectures at a shared reflexive sequence.
\item Arity invariance ensures that the number of components in a composite or chain remains fixed across reflexive and meta-reflexive cycles, preserving the ordered multiplicity of the structure even as depth increases.
\item A chain is a linear sequence of composites or relational units, ordered and connected at the same depth to propagate relational dynamics across the sequence.
\item When multiple chains and composites interconnect in non-linear or branching patterns, they assemble into a network: a structured configuration that organizes complex relational architectures through an integrated web of connections.
\item Within a network, each distinct relational unit or composite functions as a node: a positional element that serves as a hub or endpoint for relational connections.
\item The relational connections between nodes are edges: relational bindings, typically realized as tension composites, that link nodes together at a shared reflexive depth to form the network’s structural framework.
\item Node/edge role invariance ensures that within any network, the distinct functions of nodes (units) and edges (bindings) remain stable under reflexive and meta-reflexive deepening, preserving the positional and functional differentiation of network components.
\end{itemize}

\item Determining external differentiation, co-presence, and relational tension provides the minimal dynamic determination for cohering the multi-relational dynamic unit, but leaves unresolved dynamics which must be derived:

\begin{itemize}
\item What determines the coherence of the multi-relational dynamic unit?
\end{itemize}
\end{itemize}
\#\# Deriving Relational Field Coherence
\begin{itemize}
\item Given that the coherence of the multi-relational dynamic unit is determined by both its need to engage relationally with other distinct multi-relational units and its need to maintain integration within the broader relational field, it cannot not derive there are dynamics of determination:

\begin{itemize}
\item In the context of something existing, these are:

\begin{itemize}
\item The act of inter-unit relating: the act that enacts and maintains direct relational engagement between distinct multi-relational units, ensuring their differentiated coherence within shared relation.

\item The act of field-integrating: the act that enacts and maintains the coherence of the entire relational field, ensuring that the differentiated multi-relational units remain integrated within a unified relational system.
\end{itemize}
\end{itemize}
\end{itemize}


These two acts arise in tandem, each satisfying a distinct necessity:

\begin{itemize}
\item Inter-unit relating addresses local coherence between distinct units.

\item Field-integrating addresses global coherence across the entire relational field.
\end{itemize}


---

\begin{itemize}
\item Given that the acts of inter-unit relating and field-integrating are determined by how they sustain relational engagement between distinct multi-relational units and maintain overall field coherence, it cannot not derive there are dynamics of relation:

\begin{itemize}
\item In the context of something existing, these are:

\begin{itemize}
\item The condition of mutual relationality: the condition in which distinct multi-relational units remain coherently related across their boundaries while preserving differentiation.

\item The condition of field coherence: the condition in which the entire relational field is stabilized to hold all multi-relational units in integrated, coherent relation.
\item These conditions induce depth alignment: multi-relational units must align their reflexive sequences to maintain coherent engagement within the integrated field.
\end{itemize}
\end{itemize}
\end{itemize}

---

\begin{itemize}
\item Given that the conditions of mutual relationality and field coherence are determined by how they structure the interplay between differentiated units and the broader relational field, it cannot not derive there are dynamics of structure:

\begin{itemize}
\item In the context of something existing, these are:

\begin{itemize}
\item The structure of mutual relational form: the minimal structure that formalizes and stabilizes the relational dynamics between distinct multi-relational units.

\item The structure of field form: the minimal structure that formalizes and stabilizes the architecture of the entire relational field, ensuring that all differentiated units are coherently integrated.
\item These structures define closure criteria: the integrated relational field is considered closed when all participating units are coherently engaged, and no further internal differentiation is induced.
\end{itemize}
\item \textbf{\textbf{Reflexive equilibrium}} ensures that after reflexive and meta-reflexive deepening, the system stabilizes such that further reflexive acts do not alter the structural form—maintaining coherence and idempotence across all layers.
\end{itemize}
\end{itemize}

---

\begin{itemize}
\item Determining inter-unit relating, field-integrating, mutual relationality, field coherence, mutual relational form, and field form provides the minimal dynamic determination for cohering the integrated relational field unit, but leaves unresolved dynamics which must be derived:

\begin{itemize}
\item What determines the coherence of the integrated relational field unit?
\end{itemize}
\end{itemize}
\#\# Deriving Relational Meta-Boundary Coherence
\begin{itemize}
\item Given that the coherence of the integrated relational field unit is determined by the stabilized interplay of all multi-relational units within a unified field, it cannot not derive there is a dynamic of determination:

\begin{itemize}
\item In the context of something existing, this is the act of meta-boundary excluding: the act that enacts the exclusion of what lies beyond the integrated relational field—what is outside of and distinct from the entire relational field.
\end{itemize}

\item Given that the act of meta-boundary excluding is determined by how it engages with what stands outside the integrated relational field, it cannot not derive there is a dynamic of relation:

\begin{itemize}
\item In the context of something existing, this is the condition of meta-negation: the condition that maintains the distinction between the integrated relational field and what is excluded from it (its meta-outside).
\end{itemize}

\item Given that the condition of meta-negation is determined by how it structures and maintains the distinction between the integrated relational field and its outside, it cannot not derive there is a dynamic of structure:

\begin{itemize}
\item In the context of something existing, this is the structure of meta-boundary: the minimal structural form that formalizes and stabilizes the division between the integrated relational field and what is outside it.
\item This structure preserves meta-level invariants: the meta-boundary maintains positional differentiation and type-homogeneity, ensuring that meta-level dynamics remain consistent with foundational relational logic.
\end{itemize}
\item Determining meta-boundary excluding, meta-negation, and meta-boundary provides the minimal dynamic determination for cohering the meta-boundary unit, but leaves unresolved dynamics which must be derived:

\begin{itemize}
\item What determines the coherence of the meta-boundary unit?
\end{itemize}
\end{itemize}
\#\# Deriving Meta-relational Coherence
\begin{itemize}
\item Given that the coherence of the meta-boundary unit is determined by its distinction between the integrated relational field and what is outside of it, it cannot not derive there is a dynamic of determination:

\begin{itemize}
\item In the context of something existing, these are:

\begin{itemize}
\item The act of meta-reflexive relating: the act that enacts relation upon the meta-boundary itself, folding its own exclusionary act back into the relational system.

\item The act of meta-transcendence: the act that engages with what the meta-boundary posits as beyond itself, extending the relational dynamic beyond the currently defined field.
\end{itemize}
\end{itemize}
\end{itemize}


---

\begin{itemize}
\item Given that the acts of meta-reflexive relating and meta-transcendence are determined by how they engage both the meta-boundary and what is posited beyond it, it cannot not derive there are dynamics of relation:

\begin{itemize}
\item In the context of something existing, these are:

\begin{itemize}
\item The condition of meta-reflexive relation: the condition in which the meta-boundary becomes part of the relational field through its own self-engagement, stabilizing its own coherence.

\item The condition of meta-expansion: the condition in which the relational system is extended to accommodate and engage with what was previously excluded by the meta-boundary.
\end{itemize}
\end{itemize}
\end{itemize}


---

\begin{itemize}
\item Given that the conditions of meta-reflexive relation and meta-expansion are determined by how they structure the interplay between the meta-boundary, its own reflexivity, and its engagement with the beyond, it cannot not derive there are dynamics of structure:

\begin{itemize}
\item In the context of something existing, these are:

\begin{itemize}
\item The structure of meta-reflexive form: the minimal structure that formalizes and stabilizes the meta-boundary’s self-relation within the relational system.

\item The structure of meta-extension: the minimal structure that formalizes and stabilizes the expansion of the relational field to engage coherently with what was previously outside.
\item These structures induce recursive domain unfolding: closure at the meta-relational level exposes new implicit dynamics, indicating the necessity of formalizing a subsequent relational domain.
\end{itemize}
\end{itemize}
\item \textbf{\textbf{Predictive determination}} formalizes that each closed relational domain exposes implicit dynamics that compel the emergence of a new domain, enabling the architecture to forecast the necessary shape of its next formal extension.
\item Determining meta-reflexive relating, meta-transcendence, meta-reflexive relation, meta-expansion, meta-reflexive form, and meta-extension provides the minimal dynamic determination for cohering the meta-relational unit, but leaves unresolved dynamics which must be derived:
\begin{itemize}
\item What determines the coherence of the meta-relational unit?
\end{itemize}
\end{itemize}
\section*{Formalizing "Deriving Relational Dynamics" in Category-Theoretic Terms}
\label{sec:org023d4ca}
A simplicial-set representation K : Δᵒᵖ → Set
\subsection*{1. Ambient category and notation}
\label{sec:org83cc211}
\begin{itemize}
\item Work in sSet = Fun(Δᵒᵖ, Set)
\item Δⁿ : standard n-simplex
\item ∂Δⁿ : its boundary
\item Λⁿₖ : k-th horn
\item Nondegenerate n-simplex of X ∈ sSet = monomorphism Δⁿ ↪ X
\item Adjoining an n-simplex along a boundary is the pushout diagram:

\begin{verbatim}
\begin{tikzcd}
\partial\Delta^n \ar[r] \ar[d, hook] & X \ar[d] \\
\Delta^n \ar[r] & X \cup_{\partial\Delta^n}\Delta^n
\end{tikzcd}
\end{verbatim}
\end{itemize}
\subsection*{2. Base object K\textsuperscript{(0)} ("K₀")}
\label{sec:org8c19235}
Define K\textsuperscript{(0)} as the simplicial set freely generated by one nondegenerate 2-simplex
σ : Δ² ↪ K\textsuperscript{(0)}
with vertices \{δ, ¬, =\}.

Pushout description:
\begin{verbatim}
\begin{tikzcd}
\partial\Delta^2 \ar[r] \ar[d, hook] & \Delta^0\sqcup \Delta^0\sqcup \Delta^0 \ar[d] \\
\Delta^2 \ar[r, hook] & K^{(0)}
\end{tikzcd}
\end{verbatim}
\subsection*{3. Adding relation r and closing to a 3-simplex (K\textsuperscript{(1)})}
\label{sec:org89459d9}
Adjoin new vertex r and one nondegenerate 3-simplex τ : Δ³ ↪ K\textsuperscript{(1)}
with vertices δ, ¬, =, r.

Pushout:
\begin{verbatim}
\begin{tikzcd}
\partial\Delta^3 \ar[r] \ar[d, hook] & K^{(0)}\sqcup \Delta^0_r \ar[d] \\
\Delta^3 \ar[r, hook] & K^{(1)} := K^{(0)}\cup_{\partial\Delta^3}\Delta^3
\end{tikzcd}
\end{verbatim}
\subsection*{4. Reflexive triad (a, ρ, f) as 3-simplex (K\textsuperscript{(2)})}
\label{sec:org1914bd0}
Add vertices a, ρ, f and attach a nondegenerate 3-simplex κ : Δ³ ↪ K\textsuperscript{(2)}.

Pushout:
\begin{verbatim}
K^{(2)} := K^{(1)} ∪_{∂Δ³→K^{(1)}∪Δ⁰_a∪Δ⁰_ρ∪Δ⁰_f} Δ³
\end{verbatim}

Interpretation: faces of κ encode partial compositions among (a, ρ, f).
\subsection*{5. Second-order morphism r′ as a 4-simplex (K\textsuperscript{(3)})}
\label{sec:org89c205e}
Encode higher relation (a,ρ) ↦ (ρ,f) by adjoining one nondegenerate 4-simplex
θ : Δ⁴ ↪ K\textsuperscript{(3)}.

Faces of θ represent the participating 3-simplices.
\subsection*{6. Boundary formation (K\textsuperscript{(4)})}
\label{sec:org0f60ddf}
Boundary operator on chains:

\begin{verbatim}
\partial_n : C_n(K^{(3)};ℤ) → C_{n-1}(K^{(3)};ℤ)
\end{verbatim}

Generated simplicially by face maps dᵢ of K\textsuperscript{(3)}.

Represent the triplet (e, n, b) as a new simplex (2- or 3-simplex β) attached via
a boundary pushout.  The boundary subcomplex ∂K₃ is the image of
sk₂(Δ³) ↪ Δ³ ↪ K\textsuperscript{(1)} ⊂ K\textsuperscript{(4)}.
\subsection*{7. Boundary reflexivity r\textsubscript{b} : ∂K₃ → ∂K₃  (K\textsuperscript{(5)})}
\label{sec:orgb119b90}
An endomorphism of a subcomplex i : B ↪ K\textsuperscript{(4)} is a simplicial map r\textsubscript{b} : B → B.

Internalize it by adding a mapping cylinder:
\begin{verbatim}
B → Cyl(r_b) → B
\end{verbatim}

Attach Cyl(r\textsubscript{b}) along B:

\begin{verbatim}
K^{(5)} := K^{(4)} ∪_{B→Cyl(r_b)} Cyl(r_b)
\end{verbatim}
\subsection*{8. Multiplicity and differentiation (K\textsuperscript{(6)})}
\label{sec:org26e6db0}
Let 𝕀 be a small index category with objects i, morphisms δᵢⱼ : K⁵ᵢ → K⁵ⱼ, δᵢᵢ = id.
Define a functor D : 𝕀 → sSet,  i ↦ K⁵ᵢ.

Multiplicity complex:
\begin{verbatim}
K^{(6)} := colim_{i∈𝕀} K^{(5)}_i
\end{verbatim}
\subsection*{9. Field integration (K\textsuperscript{(7)})}
\label{sec:orgf2f0996}
Inter-unit relations φᵢⱼ : K⁵ᵢ ⇄ K⁵ⱼ extend 𝕀 to a span-enriched diagram 𝔇 in sSet.

Field complex and integration map:
\begin{verbatim}
F := colim 𝔇
Φ : ⨿_i K^{(5)}_i → F
K^{(7)} := F
\end{verbatim}
\subsection*{10. Meta-boundary, folding, and extension (K\textsuperscript{(8)}, K\textsuperscript{(9)})}
\label{sec:org12ca133}
Meta-boundary: monomorphism ιₘ : Bₘ ↪ F.

Meta-reflexivity and meta-transcendence:
\begin{verbatim}
Rₘ : Bₘ → F
Tₘ : Bₘ → F′
\end{verbatim}

Pushout forming meta-extension:
\begin{verbatim}
\begin{tikzcd}
B_m \ar[r,"T_m"] \ar[d,"ι_m"'] & F' \ar[d] \\
F \ar[r] & F ∪_{B_m} F'
\end{tikzcd}
\end{verbatim}

Define K\textsuperscript{(9)} := F ∪\textsubscript{Bₘ} F′.

Iterating (Tₘ ∘ Rₘ) gives chain F ⊂ F′ ⊂ F″ ⊂ ⋯.
\subsection*{11. Inductive system and colimit K\_∞}
\label{sec:orgddcf13e}
Filtered diagram:
\begin{verbatim}
K^{(0)} → K^{(1)} → K^{(2)} → ⋯
\end{verbatim}

Each iₙ is a monomorphism from corresponding pushout.

Colimit in sSet:
\begin{verbatim}
K_∞ := colim_{n∈ℕ} K^{(n)}
\end{verbatim}

Levelwise:
\begin{verbatim}
K_∞[m] = colim_n K^{(n)}[m]
\end{verbatim}

Face and degeneracy maps inherited naturally from Δ.
\subsection*{12. Optional fibrancy (inner horn filling)}
\label{sec:org5c43af0}
To ensure higher-order closure, attach fillers for all inner horns Λⁿₖ ↪ Δⁿ.

This produces a Kan complex (fibrant object) K\textsuperscript{(ω)}
and a canonical map K\textsuperscript{(n)} → K\textsuperscript{(ω)}.
\subsection*{13. Chain-level boundary (homology tie-in)}
\label{sec:org1410300}
Normalized chain complex:
\begin{verbatim}
N_*(K_∞;ℤ),  ∂ₙ = Σ_{i=0}^n (-1)^i d_i : N_n → N_{n-1}
\end{verbatim}

Encodes the "boundary formation" operator ∂.
\subsection*{14. Compact definition}
\label{sec:org22786df}
A \textbf{relational dynamics object} is a simplicial set K : Δᵒᵖ → Set
presented by a filtered diagram (K\textsuperscript{(n)}, iₙ) in sSet such that:

\begin{enumerate}
\item K\textsuperscript{(0)} is freely generated by one nondegenerate 2-simplex
with vertices \{δ, ¬, =\}.
\item For each n ≥ 0, K\textsuperscript{(n+1)} is obtained from K\textsuperscript{(n)} by
pushouts of ∂Δᵐ ↪ Δᵐ (attaching simplices for r, a, ρ, f, r′, …)
and possibly horn fillers Λᵐₖ ↪ Δᵐ for stability.
\item K := colimₙ K\textsuperscript{(n)} (i.e. K\_∞).
\end{enumerate}

This functor K : Δᵒᵖ → Set has face/degeneracy maps induced naturally
from Δ, forming a coherent simplicial hierarchy representing
the formal derivation of existence as relational dynamics.
\end{document}
